%%%%%%%%%%%%%%%%%%%%%%%%%%%%%%%%%%%%%%%%%%%%%%%%%%%%%%%%%%%%
% (01) Input Settings
%%%%%%%%%%%%%%%%%%%%%%%%%%%%%%%%%%%%%%%%%%%%%%%%%%%%%%%%%%%%
\input{settings.tex}

%%%%%%%%%%%%%%%%%%%%%%%%%%%%%%%%%%%%%%%%%%%%%%%%%%%%%%%%%%%%
% (02) DOCUMENT INFORMATION - FROM YAML METADATA
%%%%%%%%%%%%%%%%%%%%%%%%%%%%%%%%%%%%%%%%%%%%%%%%%%%%%%%%%%%%
\newcommand{\DocTitle}{شی‌گرایی: مدیریت استثنا \lr{(Exception Handling)}}

\newcommand{\CourseTitle}{برنامه‌سازی پیشرفته و کارگاه}

\newcommand{\Semester}{بهار ۱۴۰۵}

\newcommand{\FirstProfessor}{دکتر احمدیان}

\newcommand{\SecondProfessor}{}

\newcommand{\Writers}{آرمان حسینی، پریا شایان فر، پریا اصحابی}

\newcommand{\Faculty}{دانشکده ریاضی و علوم کامپیوتر}

\newcommand{\University}{دانشگاه صنعتی امیرکبیر}

%%%%%%%%%%%%%%%%%%%%%%%%%%%%%%%%%%%%%%%%%%%%%%%%%%%%%%%%%%%%
% DOCUMENT BODY
%%%%%%%%%%%%%%%%%%%%%%%%%%%%%%%%%%%%%%%%%%%%%%%%%%%%%%%%%%%%
\begin{document}
	
%%%%%%%%%%%%%%%%%%%%%%%%%%%%%%%%%%%%%%%%%%%%%%%%%%%%%%%%%%%%
% === TITLE PAGE ===
%%%%%%%%%%%%%%%%%%%%%%%%%%%%%%%%%%%%%%%%%%%%%%%%%%%%%%%%%%%%
\begin{titlepage}
	\centering
	\vspace*{1cm}
	
	%— Smaller logo (20% of text width)
	\begin{figure}[H]
		\centering
		
\includegraphics[width=0.20\textwidth]{images/logo-fa.png}
	\end{figure}
	
	\vspace{1cm}
	
	{\fontsize{18}{28}\selectfont \textbf{\CourseTitle} }
	
	\vspace{3cm}
	
	{\fontsize{22}{28}\selectfont \textbf{\DocTitle} }
	
	\vspace{3cm}
	
	{\fontsize{14}{22}\selectfont \rl{نگارش:}}\\
	\vspace{0.5cm}
	{\fontsize{16}{22}\selectfont \textbf{ \Writers } }

	
	\vspace{1cm}
	
	{\fontsize{14}{22}\selectfont \rl{استاد درس:}}\\
	\vspace{0.5cm}
	{\fontsize{16}{22}\selectfont \textbf{ \FirstProfessor } }
	
	
	\vfill
	
	{\fontsize{14}{18}\selectfont \textbf{\Faculty} }\\
	\vspace{0.5cm}
	{\fontsize{14}{18}\selectfont \textbf{\University} }
	\vspace{0.5cm}
	
	\par
	
	{\fontsize{14}{18}\selectfont \textbf{\Semester} }
	
\end{titlepage}

%— Switch to our fancy header/footer from this page onward
\clearpage

%%%%%%%%%%%%%%%%%%%%%%%%%%%%%%%%%%%%%%%%%%%%%%%%%%%%%%%%%%%%
% === TABLE OF CONTENTS ===
%%%%%%%%%%%%%%%%%%%%%%%%%%%%%%%%%%%%%%%%%%%%%%%%%%%%%%%%%%%%
\tableofcontents
\clearpage
\pagestyle{aftertitle}

%%%%%%%%%%%%%%%%%%%%%%%%%%%%%%%%%%%%%%%%%%%%%%%%%%%%%%%%%%%%
% === MAIN CONTENT ===
%%%%%%%%%%%%%%%%%%%%%%%%%%%%%%%%%%%%%%%%%%%%%%%%%%%%%%%%%%%%
\chapter{\DocTitle}
\section{مقدمه}

یکی از مفاهیم مهم در زبان جاوا، مفهوم \lr{Exception Handling} است. «
\lr{Exception} » در زبان انگلیسی به معنای «استثناء» است و شرایط استثنایی
در جاوا، شرایطی است که برنامه‌ی شما به خطا می خورد. در زبان جاوا،
\lr{Exception} (استثنا) رویدادی است که در حین اجرای برنامه
\lr{(Runtime)} رخ می‌دهد و جریان عادی اجرای دستورات را متوقف می‌کند ولی
جاوا به شما این امکان را می‌ دهد که به جای \lr{crash} کردن برنامه موقع به
وجود اومدن خطا، خودتون به شکلی دیگر با خطای به وجود آمده برخورد کنید.

خطاها همیشه نباید به \lr{crash} منجر شوند .فرض کنید برنامه‌نویس
\lr{Google} هستید؛ قاعدتاً نمی‌خواهید به محض این که یکی از کاربرها در
\lr{refresh} کردن \lr{gmail} خود به خطا خورد، تمام سرورهای \lr{gmail}
برای همه‌ی کاربرها \lr{crash} کند. این‌جاست که مدیریت خطاهای سیستم و
\lr{Exception Handling} بسیار مهم می شوند.

\section{\lr{Exception} ها}

در ادامه میخواهیم موقعیت‌هایی را بررسی کنیم که برنامه‌ی شما مممکن است با
یک \lr{Exception} مواجه شود. برنامه‌ی زیر رو توی جاوا اجرا کنید:

\begin{codebox}{java}
import java.util.Scanner;

public class Main {
    public static void main(String[] args) {  
        Scanner scanner = new Scanner(System.in);

        int[] arr = {21, 11, 2, 5, 98};

        while (true) {  
            System.out.print("Enter the index of your number (between 0 and 4): ");
            var idx = scanner.nextInt();  
  
            System.out.println("Number at that idx: " + arr[idx]);  
        }  
    }  
}
\end{codebox}

این برنامه، یه عدد به اسم ‍ از کاربر می‌گیرد و \lr{\texttt{arr[idx]}} را
خروجی می‌دهد. مثلاً، یه نمونه از اجرای این برنامه به شکل زیر است:

\begin{codebox}{text}
Enter the index of your number (between 0 and 4): 0  
Number at that idx: 21  
Enter the index of your number (between 0 and 4): 2  
Number at that idx: 2
\end{codebox}

چه اتفاقی می افتد اگر کاربر عددی خارج از بازه‌ی ۰ تا ۴ را وارد کند؟
همون‌طور که خودتان می‌دانید، به یک خطا می‌خوریم:

\begin{figure}[H]
\centering
\includegraphics[width=0.8\textwidth]{images/exception-handling-1.png}
\end{figure}

به متن این خطا توجه کنید:

\begin{codebox}{text}
Exception in thread "main" java.lang.ArrayIndexOutOfBoundsException: Index -2 out of bounds for length 5
\end{codebox}

جاوا به شما می‌گوید که یک \lr{Exception} توی \lr{thread} ای به اسم
\lr{\texttt{main}} رخ داده است. ما هنوز با مفهوم \lr{thread} آشنا
نیستیم، صرفا بدونید که برنامه‌های شما همیشه توی \lr{thread} ای به اسم
\lr{\texttt{main}} اجرا می‌شوند. این خطا همچنین به شما می‌گوید که
\lr{Exception} به وجود امده، از جنس
\lr{\texttt{ArrayIndexOutOfBoundException}} است که اسم خیلی دقیقی برای
خطای شماست! در ادامه هم یک توضیح برای علت این خطا امده، این که «ایندکس
-2، خارج از بازه‌ی ایندکس‌های معقول یک آرایه‌ی ۵تایی است».

در ادامه، حتی گفته شده که این خطا توی دقیقا چه تابعی رخ داده:

\begin{codebox}{text}
at Main.main(Main.java:13)
\end{codebox}

شما می‌توانید بفهمید که این خطا، توی متد \lr{\texttt{Main.main}} ، در خط
۱۳ام فایل \lr{\texttt{Main.java}} رخ داده است. چند نمونه مختلف از
\lr{exception} هایی که ممکن است با آن مواجه شوید عبارتند از :

\begin{itemize}
\tightlist
\item
  \lr{\texttt{IndexOutOfBoundsException}} برای زمانی که مشکلی در ایندکس
  داده شده برای لیست یا آرایه مورد نظر وجود دارد.
\item
  \lr{\texttt{NullPointerException}} به طور کل وقتی روی یک متغیر که
  \lr{null} است متدی را صدا میزنیم.
\item
  \lr{\texttt{ArithmeticException}} برای خطاهای محاسباتی مثل تقسیم بر
  صفر.
\end{itemize}

\subsection{بلوک \lr{try-catch}}

همون‌طور که توی مقدمه گفتیم، \lr{Exception} به معنای اینه که برنامه‌ی شما
به خطا خورده است. ولی خطاها، بخشی طبیعی از برنامه‌ی شما هستند. ما
نمی‌خواهیم هر بار که کاربر، بر حسب اشتباه عددی خارج از بازه‌ی ۰ تا ۴ وارد
کرد برنامه با یک خطا \lr{crash} کند و خارج شود، در عوض می‌خواهیم که به
کاربر یک پیام مناسب نشان داده شود تا عدد درستی وارد کند. پیامی مثل این:

\begin{codebox}{text}
Please enter a valid index. this index should be between 0 and 4.
\end{codebox}

برای کنترل و مدیریت خطایی که با ان مواجه شده ایم از بلوک \lr{try-catch}
استفاده میکنیم. این بلوک راه حل اصلی جاوا (و بسیاری از زبان‌های
برنامه‌نویسی دیگر) برای «مدیریت استثناها» \lr{(Exception Handling )} است
و هدف اصلی آن، جلوگیری از \lr{Crash} کردن برنامه در هنگام بروز خطا است.
این کار، با کلیدواژه‌های جدیدی به اسم \lr{\texttt{try}} و
\lr{\texttt{catch}} ممکن می شود:

\begin{codebox}{java}
try {  
    System.out.print("Enter the index of your number (between 0 and 4): ");
    var idx = scanner.nextInt();  
  
    System.out.println("Number at that idx: " + arr[idx]);  
} catch (ArrayIndexOutOfBoundsException e) {  
    System.out.println("Please enter a valid number. this number should be between 0 and 4.");  
}
\end{codebox}

کلیدواژه‌ی \lr{\texttt{try}} ، به معنای تلاش کردن است. در این کد، شما به
جاوا می‌گویید که «تلاش کن خطوط بخش \lr{\texttt{try}} را اجرا کنی، و اگر
حین اجرای این خطوط، به \lr{\texttt{Exception}} ای از جنس
\lr{\texttt{ArrayIndexOutOfBoundsException}} برخورد کردی، به جای این که
\lr{crash} کنی و از برنامه خارج شوی، به بخش \lr{\texttt{catch}} برو و
کدهای اون بخش رو اجرا کن.» دقت کنید که در پرانتز جلوی
\lr{\texttt{catch}} ، ما نوع \lr{Exception} ای که می‌خواهیم هندل کنیم را،
مثل یک تابع، توی پرانتز آورده‌ایم.

خروجی برنامه‌ی جدید شما به شکل زیر است :

\begin{codebox}{text}
Enter the index of your number (between 0 and 4): 4  
Number at that idx: 98  
Enter the index of your number (between 0 and 4): -2  
Please enter a valid number. this number should be between 0 and 4.  
Enter the index of your number (between 0 and 4): 7  
Please enter a valid number. this number should be between 0 and 4.
\end{codebox}

\subsection{\lr{catch} های متوالی}

کاربر حین اجرای برنامه‌ی شما، خطای دیگری هم میتواند ایجاد کند، این که به
جای عدد، چیزهای نامربوطی مثل اسم و رشته ورودی بدهد. برنامه را اجرا کنید
و به جای عدد، \lr{\texttt{"ali"}} را ورودی دهید تا ببینید چه اتفاقی
می‌افتد:

\begin{figure}[H]
\centering
\includegraphics[width=0.8\textwidth]{images/exception-handling-2.png}
\end{figure}

برنامه‌ی شما، یک \lr{exception} جدید داد. این بار می‌توانید ببینید که این
\lr{Exception} ، از جنس \lr{\texttt{InputMismatchException}} است. شما
می‌توانید با اضافه کردن یک \lr{\texttt{catch}} جدید به کدتان، این خطا را
هم بهتر مدیریت کنید:

\begin{codebox}{java}
try {  
    System.out.print("Enter the index of your number (between 0 and 4): ");
    var idx = scanner.nextInt();  
  
    System.out.println("Number at that idx: " + arr[idx]);  
} catch (ArrayIndexOutOfBoundsException e) {  
    System.out.println("Please enter a valid number. this number should be between 0 and 4.");  
} catch (InputMismatchException e) {  
    System.out.println("Please enter a valid number.");  
    scanner.nextLine();  
}
\end{codebox}

توی کد بالا، جاوا وقتی به خطایی برخورد میکند، اول چک می‌کند که «آیا این
خطا، از نوع \lr{\texttt{ArrayIndexOutOfBoundsException}} که توی
\lr{catch} اول امده هست یا نه؟» و اگر بود، خطوط داخل \lr{catch} اول رو
اجرا می‌کند. اما اگر نبود، به سراغ \lr{\texttt{catch}} بعدی می رود و
دوباره چک می کند که «آیا این خطا از جنس
\lr{\texttt{InputMismatchException}} هست یا نه؟» و اگر بود، خطوط
\lr{\texttt{catch}} بعدی را اجرا می‌کند.

در واقع شما ممکن است چندین \lr{\texttt{catch}} پشت سر هم داشته باشید. در
این صورت با به وجود اومدن خطا جاوا به همین ترتیب، از اول تا آخر انها را
چک می‌کند و به دستوراتِ داخل اولین \lr{\texttt{catch}} مناسب می‌رود.

خروجی برنامه‌ی جدید شما، به شکل زیر است:

\begin{codebox}{text}
Enter the index of your number (between 0 and 4): 2  
Number at that idx: 2  
Enter the index of your number (between 0 and 4): ali  
Please enter a valid number.  
Enter the index of your number (between 0 and 4): -1  
Please enter a valid number. this number should be between 0 and 4.
\end{codebox}

\section{کلاس \lr{Exception}}

کلاس \lr{\texttt{Exception}} در واقع کلاس پایه برای
\lr{Checked Exception} ها است. اگر دقت کنید، توی هر کدوم از
\lr{\texttt{catch}} ها، ما چیزی شبیه به آرگومان‌های تابع توی پرانتز ورودی
می‌گرفتیم:

\begin{codebox}{java}
catch (ArrayIndexOutOfBoundsException e) {
    // Your Code
}
\end{codebox}

چرا؟ مگر \lr{\texttt{ArrayIndexOutOfBoundsException}} کلاس است؟ آیا این
یعنی \lr{\texttt{e}} آبجکت است؟ بله!

روی \lr{\texttt{ArrayIndexOutOfBoundsException}} برید، کلیک راست کنید و
با \lr{Go To} و سپس \lr{Declaration and Usages} ، به تعریف این کلاس
برید:

\begin{figure}[H]
\centering
\includegraphics[width=0.8\textwidth]{images/exception-handling-3.png}
\end{figure}

می‌بینید که این کلاس، یک کلاس خیلی معمولی با سه‌تا کانستراکتور و یک فیلد
است که از کلاسی به اسم \lr{\texttt{IndexOutOfBoundsException}} ارث‌بری
کرده است!

توی جاوا، کلاسی به اسم کلاس \lr{\texttt{Exception}} داریم، که تمامی
کلاس‌های \lr{Exception} مختلفی که توی جاوا می‌بینید از ان ارث‌بری می‌کنند.
این کلاس، یک متد به اسم \lr{\texttt{()getMessage}} دارد \footnote{البته،
  اگر سورس کد کلاس \lr{\texttt{Exception}} را بخوانید، می‌بینید که این
  متد متعلق به کلاس ابسترکتی به اسم \lr{\texttt{Throwable}} هستند. اگر
  خواستید، راجع به \lr{\texttt{Throwable}} تحقیق کنید و بیشتر بخونید.}
که متن خطا مرتبط با \lr{Exception} فعلی‌مان را خروجی می‌دهد. در واقع زمانی
که خطایی رخ می‌دهد، جاوا یک شیء ( \lr{Object )} از نوع آن خطا می‌سازد و آن
را پرتاب ( \lr{Throw )} می‌کند. (با \lr{\texttt{throw}} کمی جلوتر آشنا
خواهید شد.)

ما می‌توانیم از این کلاس توی \lr{\texttt{catch}} هایمان استفاده کنیم. کد
روی صفحه را به کد زیر تغییر دهید:

\begin{codebox}{java}
try {  
    System.out.print("Enter the index of your number (between 0 and 4): ");
    var idx = scanner.nextInt();  
  
    System.out.println("Number at that idx: " + arr[idx]);  
} catch (Exception e) {  
    System.out.println(e.getMessage());  
}
\end{codebox}

در جاوا، بخش ( \lr{\texttt{Exception e}} ) که داخل پرانتز جلوی
\lr{\texttt{catch}} نوشته می‌شود، تعیین می‌کند که کدام نوع استثنا
\lr{(Exception)} قرار است توسط آن بلوک \lr{\texttt{catch}} مدیریت شود و
همچنین نام متغیری که استثنا به آن اختصاص می‌یابد. در اینجا اگر خطایی توی
بلوک \lr{\texttt{try}} اتفاق افتاد، جاوا به اولین \lr{\texttt{catch}}
تان نگاه می‌کند و مثل قبل از خودش می‌پرسد «آیا خطا از جنس
\lr{\texttt{Exception}} است؟» و از اون‌جایی که همه‌ی خطاها از
\lr{\texttt{Exception}} ارث‌بری کرده اند، جواب همیشه مثبت است\footnote{درواقع
  جواب همیشه‌ی همیشه مثبت نیست. نوع دیگری از خطاها وجود دارند که کلاس
  پدرشان به جای \lr{\texttt{Exception}} ، کلاس کاملا متفاوتی به اسم
  \lr{\texttt{Error}} است. ما در این‌جا به آن‌ها نمی‌پردازیم ولی اگر
  خواستید، خودتان راجع بهشون بخوانید.} و کد داخل این \lr{\texttt{catch}}
اجرا می‌شود.

توی این \lr{\texttt{catch}} ، جاوا صرفا \lr{\texttt{e.getMessage()}} ،
که متن خطایمان هست را چاپ می‌کند. می‌تونید یک بار این کد را اجرا کنید تا
خروجی ان را ببینید:

\begin{codebox}{text}
Enter the index of your number (between 0 and 4): 1  
Number at that idx: 11  
Enter the index of your number (between 0 and 4): -1  
Index -1 out of bounds for length 5  
Enter the index of your number (between 0 and 4): 7  
Index 7 out of bounds for length 5
\end{codebox}

می‌بینید که متن خطایی که قبلا دیده بودیم، یعنی
\lr{\texttt{"Index -1 out of bounds for length 5"}} ، الآن روی صفحه چاپ
می شود و برنامه \lr{crash} نمی کند. هر خطای دیگه‌ای هم رخ بدهد، دقیقا به
همین \lr{\texttt{catch}} می‌رسد و متن خطایش چاپ می‌شود.

\section{\lr{throw}}

استفاده از \lr{\texttt{throw}} در جاوا معمولاً در موقعیت‌هایی است که
می‌خواهیم به صورت دستی یک \lr{Exception} ایجاد و پرتاب کنیم تا خطایی را
که در شرایط خاصی رخ می‌دهد به شکل قابل مدیریت اطلاع دهیم. در واقع با
استفاده از این کلید واژه خطا را تولید و می‌گوییم «اگر فلان اتفاق افتاد
این خطا را با پیام مشخصی که خودمان تعیین میکنیم نمایش بده». به تفاوت
بلوک \lr{try-catch} و \lr{throw} دقت کنید. وقتی کدی داریم که ممکن است
خطا بدهد، آن را داخل \lr{\texttt{try}} می‌گذاریم.اگر خطا رخ بدهد، کنترل
برنامه به \lr{\texttt{catch}} منتقل می‌شود تا جلوگیری از \lr{crash} شدن
برنامه جلوگیری شده و با مدیریت خطا و نمایش پیام مناسب خطا رفع شود. ولی
\lr{\texttt{throw}} برای این است که خودمان یک \lr{exception} را پرتاب
کنیم.اینجا \lr{exception} را سیستم تولید نکرده، بلکه خود برنامه‌نویس با
\lr{\texttt{throw}} آن را ایجاد کرده است.

کلاس زیر را، پایین کلاس \lr{\texttt{Main}} ایجاد کنید:

\begin{codebox}{java}
class Human {
    private int age;
    
    public int getAge() {
        return age;  
    }
    
    public void setAge(int age) {
        this.age = age;  
    }  
}
\end{codebox}

سپس، کد متد \lr{\texttt{main}} را به کد زیر تغییر بدهید:

\begin{codebox}{java}
public class Main {
    public static void main(String[] args) {  
        Human ali = new Human();  
        ali.setAge(10);  
          
        System.out.println("Ali's age: " + ali.getAge());  
    }  
}
\end{codebox}

کد خیلی ساده‌ای است و خروجی‌ ان هم به صورت زیر می باشد:

\begin{codebox}{text}
ali's age: 10
\end{codebox}

ولی خب، کد زیر را نگاه کنید. توی این کد، به علی سنِ -3 را داده ایم. آیا
این کار معقول است؟

\begin{codebox}{java}
Human ali = new Human();  
ali.setAge(-3);
\end{codebox}

ما می‌خواهیم که اگر برنامه‌نویس دیگری، از کلاس \lr{\texttt{Human}} استفاده
کرد و تلاش کرد سنِ کسی را منفی کند، \lr{\texttt{Exception}} ای دریافت کند
که جلویش را بگیرد. برای این کار، متد \lr{\texttt{setAge(int age)}} کلاس
\lr{\texttt{Human}} را به شکل زیر تغییر بدید:

\begin{codebox}{java}
public void setAge(int age) {
    if (age < 0) {
        var exception = new IllegalArgumentException("Age cannot be negative");
        throw exception;  
    }

    this.age = age;  
}
\end{codebox}

توی این کد، در ابتدای متد \lr{\texttt{setAge}} چک می‌کنیم که آیا سن ورودی
داده شده منفی است یا نه. اگر بود، یک \lr{\texttt{Exception}} جدید از جنس
\lr{\texttt{IllegalArgumentException}} ایجاد می‌کنیم، متن آن را با
استفاده از کانستراکتورش \lr{\texttt{"Age cannot be negative"}} می‌گذاریم
و نهایتاً آن را \lr{\texttt{throw}} می‌کنیم. با این کار، همزمان با رسیدن
به این خطا از متد خارج هم می‌شویم. حالا تلاش کنید سن \lr{\texttt{ali}} را
منفی کنید. به خطای زیر می‌خورید:

\begin{figure}[H]
\centering
\includegraphics[width=0.8\textwidth]{images/exception-handling-4.png}
\end{figure}

متن خطایی که گرفتید را بخوانید. جاوا به شما می‌گوید که توی \lr{thread} ای
به اسم \lr{\texttt{main}} ، خطایی از جنس
\lr{\texttt{IllegalArgumentException}} گرفته است که متنش،
\lr{\texttt{"Age cannot be negative"}} است. همچنین به شما می‌گوید که این
خطا، توی متد \lr{\texttt{Human.setAge}} رخ داده که خط ۲۲ ام فایل
\lr{\texttt{Main.java}} است (این خط، خطی است که توش
\lr{\texttt{exception}} مان را \lr{\texttt{throw}} کرده بودیم).

حالا این خطا را \lr{\texttt{catch}} کنید(درواقع بعد از رسیدن به بخش
\lr{\texttt{throw}} به جای ادامه اجرای عادی خط های بعدی به بخش مدیریت
خطا میرویم):

\begin{codebox}{java}
try {  
    ali.setAge(-3);  
} catch (IllegalArgumentException e) {  
    System.out.println(e.getMessage());  
}
\end{codebox}

برنامه را اجرا کنید و خروجی جدید را ببینید:

\begin{codebox}{text}
Age cannot be negative  
ali's age: 0
\end{codebox}

می‌بینید که به جای یک خطای عجیب غریب که باعث \lr{crash} برنامه‌تان می‌شد،
فقط متن خطا روی صفحه چاپ شده و بعد از ان، سن علی روی صفحه چاپ شده است
(چرا سن علی صفر است؟). آبجکت \lr{\texttt{e}} ، که یک جورایی شبیه ورودی
\lr{\texttt{catch}} تان است، دقیقا همان آبجکت \lr{\texttt{exception}} ای
که توی متد \lr{\texttt{setAge}} ساختید و ان را \lr{\texttt{throw}}
کردید:

\begin{codebox}{java}
if (age < 0) {
    var exception = new IllegalArgumentException("Age cannot be negative");
    throw exception;  
}
\end{codebox}

البته که شما می توانستیدا کدی معادل کد بالا به شکلی کوتاه‌تر بنویسید:

\begin{codebox}{java}
if (age < 0) {
    throw new IllegalArgumentException("Age cannot be negative");  
}
\end{codebox}

این مدل دوم کد‌نویسی، استانداردتر است. از دومی در کدهای خودتان استفاده
کنید، ما اولی را نوشتیم تا توجه‌تان را به این جلب کنیم که
\lr{\texttt{Exception}} ها، در واقع همگی \lr{object} هایی معمولی‌اند.

\subsection{انواع \lr{exception} ها و کلیدواژه‌ی \texttt{throws}}

حالا که کلی با \lr{Exception} ها کار کردید، لازم است بدانید که خود
\lr{Exception} ها، دو دسته دارند: \#\#\#\# \lr{Checked Exceptions}

دسته‌ی اول از \lr{exception} ها، \lr{checked exception} ها هستند. این
خطاها انقدر مهم هستند که نمی‌توانید نادیده‌شان بگیرید! این نوع
\lr{Exception} ها در زمان کامپایل بررسی می‌شوند.یعنی کامپایلر مجبور می‌کند
که برنامه‌نویس آن‌ها را مدیریت کند. اگر مدیریت نشوند، برنامه کامپایل
نمی‌شود. اگر یک متد ممکن است همچین خطایی ایجاد کند، باید حتماً توی تعریفش
این موضوع را مشخص کنید تا بقیه‌ی برنامه‌نویس‌هایی که قرار است از آن استفاده
کنند، بدانند ممکن است همچین خطای مهمی بگیرند. از طرفی، وقتی دارید از
متدی که \lr{checked exception} دارد استفاده می‌کنید، حتماً باید توی
\lr{try-catch} ان خطا را هندل کنید که برنامه‌تان \lr{crash} نکند!

یک مثال از این \lr{exception} ها، \lr{\texttt{IOException}}است. عبارت
\lr{IO} مخفف \lr{input/output} است و این خطا وقتی رخ می‌دهد که در ورودی
گرفتن و خروجی دادن چیزی به مشکل بخورید. مثلا اگر بخواهید فایلی را بخونید
و آن فایل وجود نداشته باشد، احتمالاً این \lr{\texttt{Exception}} را
می‌گیرید. توی کد زیر، متد \lr{\texttt{readFile}} (که صرفا جنبه‌ی آموزشی
دارد!)، این \lr{\texttt{exception}} را \lr{\texttt{throw}} می‌کند:

\begin{codebox}{java}
public class Main {
    public static void main(String[] args) {}

    public static void readFile(boolean fileExists) {
        if (!fileExists) {
            throw new IOException("File does not exist");  
        }  
    }  
}
\end{codebox}

این کد را توی \lr{IntelliJ} کپی کنید. می‌بینید که کدتان کامپایل نمی‌شود و
به خطای زیر می‌خورید:

\begin{figure}[H]
\centering
\includegraphics[width=0.8\textwidth]{images/exception-handling-5.png}
\end{figure}

علت این خطا، این است که متد \lr{\texttt{readFile}} با این که
\lr{\texttt{exception}} مهمی از نوع \lr{checked exception} را
\lr{\texttt{throw}} می‌کند، ولی هیچ‌جایی در تعریف خود متد نگفته که ممکن
است این \lr{exception} را بدهد. شما باید حتماً توی تعریف این‌جور متدها،
\lr{checked exception} هایشان را مشخص کنید:

\begin{codebox}{java}
public static void readFile(boolean fileExists) throws IOException {
    if (!fileExists) {
        throw new IOException("File does not exist");  
    }  
}
\end{codebox}

توی این کد جدید، کنار تعریف متد \lr{\texttt{readFile}} ، با کلیدواژه‌ی
\lr{\texttt{throws}} مشخص کردیم که این متد ممکن است خطایی از جنس
\lr{\texttt{IOException}} بدهد. حالا کدتان کامپایل می‌شود.

سعی کنید از متد \lr{\texttt{readFile}} توی متد \lr{\texttt{main}}
استفاده کنید:

\begin{codebox}{java}
public static void main(String[] args) {  
    readFile(false);  
}
\end{codebox}

باز هم خطای کامپایل می‌گیرید! این‌ بار، علت این است که شما از متد
\lr{\texttt{readFile}} استفاده کردید، و این متد ممکن است خطای مهمی بدهد،
ولی شما این خطا را در یک \lr{\texttt{try-catch}} هندل نکردید. شما دوتا
راه‌حل دارید که بسته به شرایط، باید یکی از ان‌ها را انتخاب کنید:

\textbf{راه حل اول}: این که در یک \lr{\texttt{try-catch}} ، خطای احتمالی
متد \lr{\texttt{readFile}} را مدیریت کنید:

\begin{codebox}{java}
public static void main(String[] args) {
    try {  
        readFile(false);  
    } catch (IOException e) {  
        System.out.println(e.getMessage());  
    }  
}
\end{codebox}

با این کار، جاوا مطمئن می‌شود که این خطا توسط شما مدیریت شده و دیگر بهتان
گیر نمی‌دهد.

\textbf{راه حل دوم}: بپذیرید که متدتان نمی‌تواند این خطا را مدیریت کند و
آن را به تعریف متد \lr{\texttt{main}} هم اضافه کنید:

\begin{codebox}{java}
public static void main(String[] args) throws IOException {  
    readFile(false);  
}
\end{codebox}

حالا اگر برنامه‌نویسی از متد \lr{\texttt{main}} استفاده کرد او هم باید
مثل شما تلاش کند تا این \lr{\texttt{IOException}} را برطرف کند!

\subsubsection{\lr{Unchecked Exceptions}}

این \lr{exception} ها،در زمان \lr{Runtime} رخ می‌دهند و بر خلاف دسته‌ی
قبل، به قدری مهم نیستند که مدیریت آن‌ها اجباری باشد.

توی مثال اول، شما دیدید که متد \lr{\texttt{scanner.nextInt()}} می‌تواند
خطایی از جنس \lr{\texttt{InputMismatchException}} بدهد، و این خطا وقتی
اتفاق می‌افتاد که کاربر به جای \lr{\texttt{int}} ، تایپ‌های دیگری مثل
\lr{\texttt{string}} یا \lr{\texttt{float}} ورودی بدهد. ولی شما تا قبل
از خواندن این داک، از جلسه‌ی اول کارگاه از متد
\lr{\texttt{scanner.nextInt()}} استفاده می‌کردید، بدون این که حتی بدانید
خطای \lr{\texttt{InputMismatchException}} چیست! چه برسد به این که توی
\lr{\texttt{try-catch}} ان را مدیریت کنید یا با \lr{\texttt{throws}} آن
را به تعریف متدتان اضافه کنید!

می‌بینید که جاوا شما را مجبور به مدیریت این \lr{exception} ها نمی‌کند.
خطاهایی مثل \lr{\texttt{ArrayIndexOutOfBoundsException}} ،
\lr{\texttt{ArithmeticException}} و کلی خطای دیگر که تا الآن صفحه‌هایتان
را قرمز می‌کردند و نمی‌دانستید چی‌ هستند، همگی از این جنس هستند.

این \lr{Exception} ها، به جای این که مستقیما از کلاس
\lr{\texttt{Exception}} ارث‌بری کنند، از کلاس دیگری به اسم
\lr{\texttt{RuntimeException}} ارث‌بری کردند:

\begin{figure}[H]
\centering
\includegraphics[width=0.8\textwidth]{images/exception-handling-6.png}
\end{figure}

می‌بینید که خود این کلاس‌ هم، از \lr{\texttt{Exception}} ارث‌بری می‌کند.

\section{ساخت \lr{Exception} های جدید}

دیدید که کدهای کلاس‌هایی مثل \lr{\texttt{ArrayIndexOutOfBoundsException}}
، \lr{\texttt{RuntimeException}} و امثال ان‌ها، کدهای نسبتا ساده‌ای‌ هستند.
عموماً فقط یک سری \lr{constructor} دارند که با استفاده از آن‌ها، بتوانید
آبجکت‌های جدیدی از آن‌ها درست کنید و آن‌ها را \lr{\texttt{throw}} کنید.
مثلاً اگر کامنت‌های کلاس \lr{\texttt{RuntimeException}} را کنار بذاریم،
فقط همچین چیزی از کد باقی می‌ماند:

\begin{codebox}{java}
public class RuntimeException extends Exception {
    @java.io.Serial
    static final long serialVersionUID = -7034897190745766939L;
    
    public RuntimeException() {
        super();  
    }
    
    public RuntimeException(String message) {
        super(message);  
    }
    
    public RuntimeException(String message, Throwable cause) {
        super(message, cause);  
    }
    
    public RuntimeException(Throwable cause) {
        super(cause);  
    }
    
    protected RuntimeException(String message, Throwable cause,
                               boolean enableSuppression,
                               boolean writableStackTrace) {
        super(message, cause, enableSuppression, writableStackTrace);  
    }  
}
\end{codebox}

پس میتوانیم به راحتی کلاس های \lr{\texttt{Exception}} خودمان را بسازیم.
برای ساختن یک \lr{\texttt{Exception}} سفارشی \lr{(Custom Exception)} در
\lr{Java} باید یک کلاس جدید بسازید که از یکی از کلاس‌های
\lr{\texttt{Exception}} ارث‌بری کند و بسته به اینکه بخواهید
\lr{Exception} شما \lr{Checked} باشد یا \lr{Unchecked} ، از کلاس متفاوتی
ارث می‌برید.

مثال قبل‌تر، که در آن به یک انسان سنِ منفی می‌دادیم را در نظر بگیرید:

\begin{codebox}{java}
class Human {
    private int age;

    public int getAge() {
        return age;  
    }

    public void setAge(int age) {
        if (age < 0) {
            throw new IllegalArgumentException("Age cannot be negative");  
        }

        this.age = age;  
    }  
}
\end{codebox}

فرض کنید می‌خواهیم این‌جا، به جای \lr{\texttt{IllegalArgumentException}} ،
از \lr{\texttt{Exception}} جدیدِ خودمون یعنی
\lr{\texttt{InvalidAgeException}} استفاده کنیم. طبق رسم‌های قراردادی
جاوا، اسم تمام کلاس‌هایی که نشان‌دهنده‌ی یک \lr{\texttt{Exception}} هستند،
باید به عبارت \lr{\texttt{Exception}} ختم شوند.

کلاس زیر را، زیر کلاس \lr{\texttt{Human}} ایجاد کنید:

\begin{codebox}{java}
class InvalidAgeException extends Exception {
    public InvalidAgeException(String message) {
        super(message);  
    }
    
    public InvalidAgeException(int age) {
        super("Invalid age: " + age);  
    }  
}
\end{codebox}

برای ساخت یک \lr{\texttt{Exception}} جدید، شما باید مثل تمامی
\lr{\texttt{Exception}} های دیگر، به شکلی مستقیم یا غیر مستقیم از کلاس
\lr{\texttt{Exception}} ارث‌بری کرده باشید. توی
\lr{\texttt{InvalidAgeException}} ، ما دوتا \lr{constructor} نوشتیم که
کار برنامه‌نویس‌های دیگر را راحت کند. هر دوی این \lr{constructor} ها، صرفا
\lr{constructor} کلاس \lr{\texttt{Exception}} را با متن خطای خاص خودشان
صدا می‌زنند. توی کانستراکتور اول، این متن خطا ورودی خود کانستراکتور است و
توی کانستراکتور دوم، صرفاً با ورودی گرفتن سنِ نامعتبر، متن خطای پیش‌فرضی که
نوشتیم به کانستراکتور کلاسِ \lr{\texttt{Exception}} داده شده است. شما
همیشه باید با صدا زدن کانستراکتورهای خود کلاس \lr{\texttt{Exception}}
متن خطا را تنظیم کنید.

بیاین تا این \lr{\texttt{Exception}} رو توی متد
\lr{\texttt{setAge(int age)}} استفاده کنیم:

\begin{codebox}{java}
class Human {
    private int age;

    public int getAge() {
        return age;  
    }

    public void setAge(int age) {
        if (age < 0) {
            throw new InvalidAgeException(age);  
        }

        this.age = age;  
    }  
}
\end{codebox}

توی متد \lr{\texttt{setAge}} ، این بار به جای
\lr{\texttt{IllegalArgumentException}} از \lr{\texttt{Exception}}
اختصاصی خودمان استفاده کردیم. اما یک خطای کوچک داریم:

\begin{figure}[H]
\centering
\includegraphics[width=0.8\textwidth]{images/exception-handling-7.png}
\end{figure}

این خطا، شبیه همان خطایی‌‍ است که توی بخش \lr{checked exception} ها بهش
برخوردیم. \lr{Exception} جدیدی که شما نوشتید، از جنس
\lr{checked exception} است. برای این که این‌طور نباشد، کافیست مثل
\lr{unchecked exception} ها از \lr{\texttt{RuntimeException}} ارث‌بری
کنید:

\begin{codebox}{java}
class InvalidAgeException extends RuntimeException {
    // CODE CODE CODE
}
\end{codebox}

با این کار، خطای خود را برطرف و کدتان کامپایل می‌شود. کد زیر را برای
\lr{\texttt{catch}} کردن این خطا توی متد \lr{\texttt{main}} بنویسید:

\begin{codebox}{java}
var ali = new Human();

try {  
    ali.setAge(-7);  
} catch (InvalidAgeException e) {  
    System.out.println("Exception: " + e.getMessage());  
}
\end{codebox}

حالا کدتان را اجرا کنید:

\begin{codebox}{text}
Exception: Invalid age: -7
\end{codebox}

شما \lr{\texttt{exception}} جدید خود را ساختید! بیاید به طور خلاصه به
مراحلی که برای ساخت \lr{exception} جدید خود طی کردید نگاهی بیندازیم :

\begin{enumerate}
\def\labelenumi{\arabic{enumi}.}
\tightlist
\item
  یک کلاس جدید ساختید.
\item
  از \lr{\texttt{Exception}} یا \lr{\texttt{RuntimeException}} ارث‌بری
  کردید. ( \lr{checked} یا \lr{unchecked} بودن خطای خود را با توجه به
  نوع ارث‌بری تعیین کردید.)
\item
  یک \lr{constructor} تعریف کرید تا بتوانید پیام خطا و علت خطا را منتقل
  کنید.
\item
  با \lr{\texttt{throw}} آن را پرتاب کردید.
\item
  اگر \lr{Checked} بود، با \lr{try-catch} یا \lr{throws} مدیریت کردید.
\end{enumerate}

\section{چه چیزی یاد گرفتیم}

توی این داک فهمیدیم که:

\begin{itemize}
\tightlist
\item
  بلوک‌های \lr{\texttt{try-catch}} دقیقا چه چیزی هستند و چطور می‌شود از
  آن‌ها استفاده کرد.
\item
  خود \lr{exception} چیست؟ انواع مختلف آن چه هستند و کلاس
  \lr{\texttt{Exception}} چه چیزهایی توی خودش نگه داشته است.
\item
  چطور می‌شود \lr{exception} های جدید خودمان را درست کنیم و از آن‌ها
  استفاده کنیم.
\item
  کلیدواژه‌ی \lr{\texttt{throw}} چیست و چطور می‌توانیم در یک برنامه به
  برنامه‌‌نویس‌های دیگر خطا بدهیم.
\end{itemize}

\section{منابع بیشتر}

اگر دوست داشتید، میتوانید
\href{https://docs.oracle.com/javase/tutorial/essential/exceptions/index.html}{داک
رسمی \lr{Oracle} برای \lr{Exception} ها} را مطالعه کنید. این داک‌ها
پیشرفته بوده و میتوانید کلی مطالب جدید از انها یاد بگیرید.

همچنین، راجع به چیزهایی مثل \lr{\texttt{Error}} و
\lr{\texttt{Throwable}} توی جاوا تحقیق کنید. هر دوی اون‌ها مباحثی پیشرفته‌
هستند که یادگیری شما رو تکمیل می‌کنند. برای مطالعه‌ی اون‌ها، علاوه بر
اینترنت می‌توانید از فصل «Error Handling» کتاب «
\lr{Learning Java - An Introduction to Real-World Programming with}
Java» هم استفاده کنید:

\begin{figure}[H]
\centering
\includegraphics[width=0.8\textwidth]{images/exception-handling-8.png}
\end{figure}
	
	
	
	
%%%%%%%%%%%%%%%%%%%%%%%%%%%%%%%%%%%%%%%%%%%%%%%%%%%%%%%%%%%%
\end{document}